\documentclass[11pt,a4paper]{article}
\usepackage[utf8]{inputenc}
\usepackage[T1]{fontenc}
\usepackage{amsmath,amssymb}
\usepackage{booktabs}
\usepackage{graphicx}
\usepackage{hyperref}
\usepackage[margin=1in]{geometry}
\usepackage{natbib}
\usepackage{xcolor}
\usepackage{float}
\usepackage{caption}
\usepackage{subcaption}

\hypersetup{colorlinks=true,linkcolor=blue!60!black,citecolor=green!50!black,urlcolor=blue!70!black}

\title{Early-Layer Locality Is a Training Curriculum,\\Not an Architectural Requirement}

\author{
Nicol\'as Veraz\\
Independent Researcher\\
\texttt{github.com/nicoveraz/neurogen}
}

\date{}

\begin{document}
\maketitle

\begin{abstract}
That early transformer layers should attend locally while late layers attend globally is well established: models learn this profile when attention span is made learnable \citep{sukhbaatar2019adaptive}, ablations confirm lower layers do not need long range \citep{rae2020deep}, schedules that ramp window size shallow-to-deep improve quality \citep{xu2025mswa}, and production models ship local--global hybrids \citep{jiang2023mistral,gemma3}. \textbf{We do not claim this result.} We ask instead \emph{when} the constraint is needed, and find that it is needed only early. Training a 3.4M-parameter transformer with a quartic depth-wise window schedule for the first 10k of 20k steps and then switching to ordinary full attention retains the entire benefit: $+1.05$\% and $+1.95$\% val\_bpb over a paired full-attention baseline on the two seeds that completed, versus $+1.26$\% and $+1.40$\% for keeping the windows throughout. Whether removal is \emph{better} than keeping them is not resolved---the two seeds disagree in sign---but removal costs nothing, which is the claim with a practical consequence: the locality prior belongs in the training recipe, not in the deployed architecture. Two measurements show what the curriculum leaves behind. Attention entropy at layer~0 remains 52\% below baseline after 100k steps, five times longer than the window phase, so the specialization is a lasting structural feature rather than a transient effect of the mask. Gradient covariance effective rank drops from 48.6 at full attention to 17.2 at window~8, with 96.7\% of gradient variance in the top component versus 6.1\%---the constraint forces early layers into a coherent, low-dimensional update subspace. A per-layer ablation localizes the effect to layers~0--1 and shows late-layer locality actively hurts, reproducing the known result in our setting. Four further experiments eliminate gradient-noise removal, softmax coupling, and landscape smoothness as mechanisms; implicit regularization remains inconclusive at $n=2$. At 125M parameters on FineWeb-Edu the windowed schedule gives $+2.6$\% at 20k steps across 5 paired seeds, with one seed negative (paired-$t$ $p=0.045$, exact paired permutation $p=0.063$). The 3.4M removal result rests on two seeds, and we show that the endpoint measurement itself carries a standard deviation of ${\approx}0.007$ bpb against a 0.0136 effect---which the paired design cancels, but only down to a residual of 0.0023 for the removal arm. We therefore pre-register five experiments that would settle the claim: the five-seed replication, a diagnosis of the divergence seen at the switch, a switch-point sweep, a fixed evaluation set, and a converged 125M three-arm comparison.
\end{abstract}

\section{Introduction}

Restricting early transformer layers to a local attention window and letting later layers attend globally is, by now, ordinary practice. It shows up as a learned property when span is made trainable \citep{sukhbaatar2019adaptive}, as an ablation result \citep{rae2020deep}, as an explicit depth-wise schedule \citep{xu2025mswa}, and as a shipped architecture \citep{jiang2023mistral,gemma3}. In every one of these treatments the locality is a property \emph{of the model}: it is present at initialization, present throughout training, and present at inference.

This paper asks a different question. If the local-to-global profile is useful, is it useful because the trained network needs it, or because the \emph{training process} needs it? These have different consequences. If the network needs it, the constraint is architectural and must be shipped. If only training needs it, the constraint is a curriculum---something to apply early and then discard, leaving a standard architecture at inference.

We test this directly by removing the windows partway through training. A 3.4M-parameter transformer trained with a quartic depth-wise window schedule for the first half of its step budget, then switched to ordinary full attention for the second half, ends up at least as good as one that kept the windows for the whole run, and better than a full-attention baseline on both seeds that completed. The benefit does not depend on the constraint remaining in place.

Two measurements indicate why. Attention entropy in the early layers stays far below baseline at $5\times$ longer training than the window phase, so the structure the curriculum induces is not maintained by the mask---it persists once the mask is gone. And under a restricted window the gradient covariance of the early-layer projections collapses to a low effective rank, so the constraint acts by making early updates coherent rather than by removing noise. A per-layer ablation shows the effect lives in layers~0--1 and that late-layer locality is harmful, which reproduces \citet{rae2020deep} in our setting and localizes it.

\paragraph{Motivation (not evidence).} This work began from an analogy to cortical critical periods, in which constrained connectivity during early development shapes functional architecture that persists after the constraints relax. The analogy is what suggested testing removal rather than presence. It is not tested by any experiment here, and no result below depends on it; we state it once and do not appeal to it again.

\paragraph{Contributions.}
\begin{itemize}
    \item \textbf{The locality constraint is removable.} Windows applied for the first 10k of 20k steps and then dropped retain the full benefit over a paired baseline ($+1.05$\%, $+1.95$\% on two seeds). This reframes a static architectural prior as a transient training curriculum (\S\ref{sec:removal}).
    \item \textbf{Two measurements of what persists.} Early-layer attention entropy stays 52\% below baseline at 100k steps; gradient covariance effective rank drops $48.6 \rightarrow 17.2$ under a window-8 constraint (\S\ref{sec:mechanism}).
    \item \textbf{A localization of the known result.} Constraining layers~0--1 alone matches or exceeds the full schedule; constraining only the last layer is worse than baseline (\S\ref{sec:ablation}).
    \item \textbf{An honest scale probe and a pre-registration.} $+2.6$\% at 125M across 5 paired seeds with one negative, and a stated protocol for the five experiments that would make the curriculum claim publishable (\S\ref{sec:scale}, \S\ref{sec:limitations}).
\end{itemize}

\section{Related Work}

\paragraph{Layer-wise attention range is an established prior.} \citet{sukhbaatar2019adaptive} make each attention head's span a learned parameter and find that trained models settle on short spans in lower layers and long spans in upper layers---the profile our schedule imposes is one gradient descent discovers on its own. \citet{rae2020deep} reach the same conclusion by ablation, showing lower layers can be restricted to a short range with no loss. \citet{xu2025mswa} impose it explicitly: MSWA varies window size across both heads and layers and progressively increases window allocation from shallow to deep layers, reporting gains in both quality and efficiency. Our per-layer ablation (\S\ref{sec:ablation}) reproduces this and adds nothing to it. We cite it as the setting our result operates in, not as our finding.

\paragraph{Local--global hybrids in production.} \citet{jiang2023mistral} use sliding-window attention in Mistral~7B to handle arbitrary-length sequences at reduced inference cost. \citet{gemma3} increase the ratio of local to global attention layers and deliberately keep the local span short, as a memory-management strategy for long contexts. Both ship the constraint permanently---at training time and at inference. Our result says the training-time half is what matters, and that the inference-time half can be dropped in settings where the goal is model quality rather than long-context cost.

\paragraph{Curriculum precedents on adjacent axes.} The closest prior work is Shortformer \citep{press2021shortformer}, which trains first on short subsequences and then on longer ones, improving perplexity and cutting training time by 1.65$\times$. That is a curriculum over \emph{how much context exists}; ours is a curriculum over \emph{how much of the available context each layer may use}, at fixed sequence length. SWAT \citep{fu2025swat} trains \emph{with} sliding-window attention specifically to close the gap between windowed inference and full-attention training, retaining windows at inference---the opposite of the deployment recommendation here. \citet{cabannes2025short} vary window size stochastically during training and find it outperforms a fixed window, but again keep the hybrid at inference and target long-context memorization rather than a discardable warmup. To our knowledge no prior work removes a depth-wise locality constraint partway through training and reports that the benefit survives.

\paragraph{Positioning.} We do not claim the architecture; we claim the mechanism. That early layers should attend locally and late layers globally is established. Our contribution is that this locality is a property of the training trajectory rather than of the trained model: the constraint can be removed partway through training with no loss, so it belongs in the training recipe, not in the deployed architecture.

\paragraph{Other background.} Sliding window attention \citep{beltagy2020longformer} and FlashAttention-2 \citep{dao2023flashattention2} supply the efficient implementation our 125M runs use. Structured-initialization work---mimetic initialization \citep{trockman2023mimetic}, T-Fixup \citep{huang2020tfixup}, $\mu$P \citep{yang2022mup}---pursues a related goal by constructing weights analytically rather than constraining computation. Induction heads \citep{olsson2022incontext} motivated the pre-wiring variants reported in Appendix~\ref{app:search}.

\section{Method}

\subsection{Models and data}

\textbf{3.4M model:} GPT-style transformer, depth 4, channels 256, 4 heads, 256-token context. Training data: TinyStories \citep{eldan2023tinystories}. Hardware: Apple M1 Pro (MPS). Metric: validation bits-per-byte (val\_bpb). Throughput is identical across architectures at 4.8 steps/sec.

\textbf{125M model:} GPT-2 Small architecture (12 layers, 768 dim, 12 heads, 1024 context). Training data: FineWeb-Edu \citep{penedo2024fineweb}. Hardware: NVIDIA H100 80GB. RoPE, RMSNorm, bfloat16 mixed precision, and FlashAttention-2 with native sliding-window support. Windowed runs are slightly \emph{faster} than baseline (2.83--2.84 vs 2.73--2.78 steps/sec) because the sliding window computes fewer attention scores.

\subsection{Depth-wise window schedule}

Each layer $l$ (of $L$) has a maximum attention window
\begin{equation}
    w(l) = \min\!\Big(T,\; b + \Big(\frac{l+1}{L}\Big)^{\!\gamma} (T - b)\Big),
    \label{eq:window}
\end{equation}
where $T$ is the sequence length, $b$ is a base floor (8 for the 3.4M model, 16 for 125M), and $\gamma$ is the growth exponent. Position $i$ attends to $[\max(0, i - w(l) + 1),\, i]$. For $\gamma = 4$ at depth 4 this gives windows $[8, 23, 86, 256]$. We searched exponents 0.5--12.0 plus sigmoid, logarithmic, exponential, and Fibonacci schedules; $\gamma \in [3,4]$ performed best at depth 4.

For the per-layer ablation we also support explicit window lists (e.g.\ \texttt{list:8,23,256,256}), which bypass Eq.~\ref{eq:window} and take the listed widths verbatim. All callers at both scales share one schedule implementation, so the 3.4M and 125M operators differ only in the base floor $b$ and the mask dtype.

\subsection{Mid-training removal protocol}
\label{sec:protocol}

The removal experiment trains three arms for 20{,}000 steps each at a fixed seed:
\begin{itemize}
    \item \textbf{A}: full attention throughout.
    \item \textbf{B}: quartic windows throughout.
    \item \textbf{F}: quartic windows for steps 0--9{,}999, then full attention for steps 10{,}000--20{,}000.
\end{itemize}
The learning-rate schedule is a single cosine over the full 20k steps with a 200-step warmup, identical across arms and \emph{continuous through the switch}---no term in the schedule is keyed to the switch point. Only the attention mask changes at step 10{,}000. All three arms share the seed, hence the same initialization and the same data order.

\subsection{Paired-seed statistics}
\label{sec:stats}

Every variant shares an initialization with its baseline at the same seed, so the runs are paired and per-seed differences are the unit of analysis. We report the exact one-sided sign-flip permutation test over the paired differences (which floors at $1/2^n$: $1/32 = 0.031$ at $n=5$, $1/4 = 0.25$ at $n=2$), the two-sided paired $t$-test, and the paired effect size $d_z$. We never use an unpaired test on these data.

This matters because the effect is smaller than the hardware noise floor. Run-to-run variance on MPS for otherwise-identical configurations is ${\approx}0.055$ bpb, while the paired residual standard deviation for the quartic variant is ${\approx}0.004$ bpb---roughly $14\times$ smaller. An unpaired comparison at these sample sizes would be uninformative; the paired one is not. Throughout, we report the number of positive paired differences alongside every mean, and any diverged run is reported rather than dropped from the denominator.

\section{Result 1: The locality constraint can be removed mid-training}
\label{sec:removal}

Table~\ref{tab:curriculum} gives the per-seed outcome of the protocol in \S\ref{sec:protocol}.

\begin{table}[H]
\centering
\caption{Removing windows at the halfway point (20k steps, 3.4M, matched init and data order per seed). $\Delta$ columns are relative val\_bpb improvement over that seed's own arm~A.}
\label{tab:curriculum}
\begin{tabular}{rcccccc}
\toprule
Seed & A: full & B: quartic & F: quartic$\rightarrow$full & B vs A & F vs A & F vs B \\
\midrule
42 & 0.9041 & 0.8927 & 0.8946 & $+1.26$\% & $+1.05$\% & $-0.22$\% \\
137 & 0.8913 & 0.8788 & 0.8739 & $+1.40$\% & $+1.95$\% & $+0.56$\% \\
256 & 0.8941 & 0.8830 & \textit{diverged (NaN)} & $+1.24$\% & --- & --- \\
\midrule
Mean (42, 137) & 0.8977 & 0.8857 & \textbf{0.8842} & $+1.33$\% & $\mathbf{+1.50}$\% & $+0.17$\% \\
\bottomrule
\end{tabular}
\end{table}

\paragraph{What the data support.} Arm~F beats arm~A on both completed seeds ($2/2$ positive paired differences). The benefit of the locality constraint therefore survives its removal: the windows are not doing ongoing work in the second half of training. This is the claim the paper rests on.

\paragraph{What the data do not support.} Whether F is \emph{better} than B is unresolved. The two seeds disagree in sign: on seed 42, F is $0.22$\% worse than B; on seed 137, F is $0.56$\% better. The mean $+0.17$\% is an average across opposite signs, and the exact sign-flip permutation over the two paired differences returns $p = 0.50$. We therefore report ``removal costs nothing,'' not ``removal helps.'' An earlier draft of this work reported the $+1.50$\% vs $+1.33$\% means without the per-seed signs and read them as evidence that windows eventually become a ceiling; that reading is not supported at $n=2$.

\paragraph{Statistical power.} With two completed seeds the exact paired permutation test floors at $p = 1/4 = 0.25$, so no arrangement of two seeds can reach conventional significance. The F-vs-A result is a consistent direction at $n=2$, not a significant one. \S\ref{sec:limitations} pre-registers the five-seed replication.

\paragraph{The switch is a shock, and the divergence is stochastic.} One of the three runs in Table~\ref{tab:curriculum}---the one at seed 256---produced NaN after the switch, and is reported here rather than excluded. It would be natural to read that as a property of seed 256. It is not. In a five-seed replication (three seeds complete at the time of writing) the \emph{same} seed completed normally, reaching 0.8865 with no divergence, so the failure is a stochastic event under this recipe rather than something a particular initialization causes. Across the original and replication runs the observed rate is one divergence in six completed runs of arm~F.

What every run does take is a large transient shock. Per-step instrumentation around the switch (Table~\ref{tab:shock}) shows the pattern is nearly identical on all three replicated seeds: loss jumps ${\approx}3.3\times$, pre-clip gradient norm jumps $16$--$20\times$, and recovery to pre-switch values is complete within about 400 steps.

\begin{table}[H]
\centering
\caption{The switch shock, measured per step around step 10{,}000 (three seeds of the five-seed replication). Gradient clipping is at 1.0, so the gradient itself is bounded---the Adam update is not.}
\label{tab:shock}
\begin{tabular}{lcccc}
\toprule
 & pre-switch & at switch & peak, first 100 steps & by step 10{,}400 \\
\midrule
Loss & 0.68 & 2.21--2.37 & --- & 0.67 \\
Gradient norm (pre-clip) & 0.136 & 2.13--2.72 & --- & 0.136 \\
Max Adam update & $3.5\times10^{-3}$ & $3.6\times10^{-3}$ & $7.5\times10^{-3}$ & $3.3\times10^{-3}$ \\
\bottomrule
\end{tabular}
\end{table}

The last row is the informative one. Gradient-norm clipping at 1.0 bounds the gradient, and the spike to 2.1--2.7 is duly clipped---but the largest Adam update still climbs $2.1\times$ over the following ${\sim}100$ steps. That is the signature of a stale second moment: $v$ was accumulated under window-8 gradients, so $m/\sqrt{v}$ is large for parameters whose historical $v$ is small, and clipping the gradient does not bound it. Of the three candidate causes we posited---(i) the attention implementation changes at the switch, from an explicit masked softmax to a fused causal kernel; (ii) a stale Adam second moment; (iii) the layer-0 softmax denominator jumping from 8 terms to up to 256 in one step---the trace is direct evidence for (ii) as the amplifier, while (i) and (iii) remain candidates for the initial spike. The learning-rate schedule is ruled out by inspection: a single cosine over the full horizon, with no discontinuity at the switch. \S\ref{sec:limitations} specifies the controls that separate the remaining causes. We do not present the curriculum as free of this: the naive hard switch is a shock that usually recovers and occasionally does not.

\section{Result 2: The active ingredient is early-layer locality}
\label{sec:ablation}

This section reproduces a known result in our setting and localizes it. \citet{sukhbaatar2019adaptive}, \citet{rae2020deep}, and \citet{xu2025mswa} all establish that lower layers want a restricted range while upper layers need global context. We confirm it here in order to identify which part of the schedule the curriculum in \S\ref{sec:removal} is actually delivering.

We replaced the quartic schedule with explicit per-layer window lists, training each at seed 42 for 20k steps. Because all configurations share the same initialization \emph{and} the same data order, the only difference is the window pattern (Table~\ref{tab:ablation}).

\begin{table}[H]
\centering
\caption{Per-layer window ablation (seed 42, 20k steps; matched init and data order). Windows are $[L_0,L_1,L_2,L_3]$; ``\% of gain'' is the fraction of the baseline$\rightarrow$quartic gap (0.0114 bpb) recovered.}
\label{tab:ablation}
\begin{tabular}{llccr}
\toprule
Config & Windows & Final bpb & vs Baseline & \% of gain \\
\midrule
Baseline & $[256,256,256,256]$ & 0.9041 & --- & --- \\
Only $L_0$ & $[8,256,256,256]$ & 0.8952 & $+0.98$\% & 78\% \\
Only $L_0,L_1$ & $[8,23,256,256]$ & \textbf{0.8880} & $\mathbf{+1.79}$\% & 141\% \\
Quartic & $[8,23,86,256]$ & 0.8927 & $+1.26$\% & 100\% \\
No $L_0$ & $[256,23,86,256]$ & 0.8979 & $+0.69$\% & 54\% \\
Only $L_3$ (reversed) & $[256,256,256,8]$ & 0.9127 & $-0.95$\% & $-75$\% \\
\bottomrule
\end{tabular}
\end{table}

Windowing layer~0 alone recovers 78\% of the quartic gain; windowing the first two layers and leaving $L_2, L_3$ at full attention exceeds the full schedule; and the reversed control---windowing only the last layer---is \emph{worse} than baseline. It is specifically early locality that helps, not locality per se. A practical corollary: the gradual quartic ramp is not essential. ``Early layers local, rest full'' matches or beats it, which is consistent with hybrid designs in which the exact schedule shape matters little \citep{gemma3,xu2025mswa}.

\textbf{Caveat.} This is a single seed. The large effects---early-local recovers most of the gain, the reversed control hurts---are robust to the noise floor. The modest margin by which Only~$L_0,L_1$ beats Quartic (0.0047 bpb) is not, and should be replicated before being asserted; \S\ref{sec:limitations} pre-registers the criterion.

\section{What the curriculum leaves behind}
\label{sec:mechanism}

If removing the windows preserves the benefit, something the windows created must persist without them. Two measurements characterize it.

\subsection{Attention entropy persists at $5\times$ longer training}

We measured per-layer attention entropy $H = -\sum_j \alpha_{ij} \log \alpha_{ij}$ on validation data (100 batches), averaged over heads, query positions, and batches, for baseline and quartic checkpoints at both 20k and 100k steps (Table~\ref{tab:entropy}).

\begin{table}[H]
\centering
\caption{Attention entropy per layer at 20k and 100k steps (seed 42, 100 batches).}
\label{tab:entropy}
\begin{tabular}{rcccccc}
\toprule
 & \multicolumn{3}{c}{20k steps} & \multicolumn{3}{c}{100k steps} \\
\cmidrule(lr){2-4} \cmidrule(lr){5-7}
Layer & Baseline & Quartic & Change & Baseline & Quartic & Change \\
\midrule
0 & 1.994 & 0.852 & $-57.3$\% & 1.860 & 0.898 & $-51.7$\% \\
1 & 1.506 & 0.916 & $-39.2$\% & 1.933 & 1.248 & $-35.4$\% \\
2 & 2.466 & 2.183 & $-11.5$\% & 2.722 & 2.489 & $-8.6$\% \\
3 & 2.120 & 2.449 & $+15.5$\% & 2.382 & 2.605 & $+9.4$\% \\
\bottomrule
\end{tabular}
\end{table}

At 100k steps the quartic model still shows 52\% lower entropy at layer~0 and 35\% lower at layer~1; the gap narrows only slightly from 20k. The final layer consistently shows \emph{higher} entropy, using its full span to integrate globally over the local features built below. Effective attention spans at 100k confirm the locality is learned rather than merely masked: quartic layer~0 attends to ${\sim}2$ of its 8 permitted tokens, while baseline layer~0 attends to ${\sim}8$ of 256. The structure is a property of the weights, not of the mask---which is what \S\ref{sec:removal} predicts.

\subsection{Gradient covariance rank: the constraint makes early updates coherent}

On a frozen baseline checkpoint we collected 50 gradient samples for the layer-0 Q/K weights at five window sizes and computed the SVD of the gradient matrix (Table~\ref{tab:coupling}).

\begin{table}[H]
\centering
\caption{Gradient covariance rank vs window size (layer 0, Q/K weights, 50 passes).}
\label{tab:coupling}
\begin{tabular}{rccc}
\toprule
Window & Effective rank & Var.\ top-1 & Var.\ top-5 \\
\midrule
8 & 17.2 & 96.7\% & 97.3\% \\
32 & 45.5 & 34.0\% & 45.3\% \\
64 & 48.4 & 9.2\% & 25.2\% \\
128 & 48.6 & 5.9\% & 23.0\% \\
256 & 48.6 & 6.1\% & 22.8\% \\
\bottomrule
\end{tabular}
\end{table}

Effective rank drops from 48.6 at full attention to 17.2 at window~8---a $2.8\times$ reduction---and 96.7\% of gradient variance concentrates in the top singular component versus 6.1\% at window~256. Under the constraint, each gradient step updates a coherent low-dimensional subspace rather than a diffuse high-dimensional one. This is a plausible route by which an early constraint fixes a hierarchy that then persists, though we have not shown the causal link directly.

A complementary measurement points the same way. On a frozen trained checkpoint, gradient noise norm is approximately constant across window sizes (0.0052--0.0058), while signal norm rises $18\times$ from window~256 (0.0017) to window~8 (0.0323). Windows increase gradient \emph{coherence}; they do not reduce noise.

\subsection{Hypotheses eliminated, and one that is not}

Table~\ref{tab:eliminated} summarizes four mechanism experiments that rule alternatives out, plus one that does not resolve.

\begin{table}[H]
\centering
\caption{Mechanism hypotheses tested. Experiment numbering follows the repository.}
\label{tab:eliminated}
\begin{tabular}{p{0.26\textwidth}p{0.42\textwidth}l}
\toprule
Hypothesis & Evidence & Verdict \\
\midrule
Gradient noise removal (Exp 1) & Noise norm constant 0.0052--0.0058 across windows 8--256; only signal changes & Eliminated \\
Softmax coupling contamination (Exp 2) & Noise fraction 4.0--6.9\% across all 4 layers & Eliminated \\
Variance reduction (Exp 3) & Larger batch cannot replicate the effect at equal tokens & Eliminated ($n=1$) \\
Landscape smoothness (Exp 6) & Quartic-trained models have \emph{lower} gradient stability on both seeds (0.0335 vs 0.0358 at seed 42; 0.1101 vs 0.1121 at seed 137) & Eliminated \\
Ongoing structural constraint (Exp 7) & Removing windows preserves the benefit (\S\ref{sec:removal}) & Eliminated \\
\midrule
Implicit regularization (Exp 4) & Seed 42: quartic gap \emph{larger} ($+2.48$\% vs $+0.32$\%). Seed 137: quartic gap \emph{smaller} ($-2.70$\% vs $+0.80$\%). Seeds contradict & \textbf{Inconclusive} at $n=2$ \\
\bottomrule
\end{tabular}
\end{table}

\paragraph{On the variance-reduction control (Exp 3).} If windows worked by reducing gradient variance, larger batches should replicate them. At 2000 optimizer steps, seed 42 only: full attention at batch 32 reaches 1.242 bpb (16.4M tokens), quartic windows at batch 32 reach 1.224 ($+1.45$\%, same tokens), and full attention at batch 128 reaches 1.086 but on $4\times$ the data (65.5M tokens). Compared at equal tokens (16M), the batch-128 run is at 1.357---worse than either batch-32 run. Larger batches look better only because they saw more data; at equal token budget the windowed run wins and the larger-batch control loses. This is a single seed and we report it as such. A batch-256 arm was launched but did not complete its step budget and is excluded.

\paragraph{On implicit regularization (Exp 4).} We measured train and validation bpb on frozen checkpoints to compute the generalization gap. The two available seeds disagree in direction. At $n=2$ with contradictory signs this experiment neither rules implicit regularization in nor out, and we do not count it among the eliminations. An earlier draft reported only the seed-42 direction and claimed a clean elimination.

\section{Scale: 125M parameters}
\label{sec:scale}

Table~\ref{tab:125m_seeds} gives the per-seed gap at a matched 20k steps across all five seeds. Each seed's baseline and quartic arms share an LR schedule, so the within-seed gap is valid even though seeds 42/137 run a 50k schedule read at step 20k while 256/789/1337 are dedicated 20k runs.

\begin{table}[H]
\centering
\caption{125M per-seed gap at a matched 20k steps (all 5 seeds, paired).}
\label{tab:125m_seeds}
\begin{tabular}{rccc}
\toprule
Seed & Baseline bpb & Quartic bpb & Gap \\
\midrule
42 & 4.101 & 3.894 & $+5.04$\% \\
137 & 4.090 & 3.989 & $+2.47$\% \\
256 & 3.696 & 3.704 & $\mathbf{-0.21}$\% \\
789 & 3.578 & 3.503 & $+2.09$\% \\
1337 & 3.620 & 3.492 & $+3.53$\% \\
\midrule
Mean & --- & --- & $\mathbf{+2.6}$\% (sd 1.94) \\
\bottomrule
\end{tabular}
\end{table}

Four of five seeds are positive; seed 256 is negative. The exact one-sided sign-flip permutation test over the five paired differences gives $p = 2/32 = 0.063$, the paired $t$-test $p = 0.045$, and $d_z = 1.29$. This is a weaker result than the 3.4M one (where all 5 seeds were positive and the permutation test hit its floor of 0.031), and we report it as a suggestive scale probe rather than a demonstrated scaling law.

\paragraph{Convergence.} Neither 50k-step arm is converged. Over the final 10k steps the baselines are still descending at 0.0098 (seed 42) and 0.0084 (seed 137) bpb per 1k steps, and the quartic arms are descending \emph{faster} still (0.0208 and 0.0166). A larger gap at 50k is therefore what an unconverged pair of curves produces regardless of whether the asymptotic gap is larger, so we do not read the 50k numbers as a scaling result. They are reported in Appendix~\ref{app:125m50k} with that label. Determining whether the gap is real requires a single fully-annealed horizon; the protocol is in \S\ref{sec:limitations}.

\paragraph{Throughput.} Windowed 125M runs are slightly faster than baseline (2.83--2.84 vs 2.73--2.78 steps/sec), because FlashAttention's sliding window computes fewer attention scores. Under the curriculum recipe this advantage applies only to the windowed phase, and inference runs at standard full-attention cost.

\section{Limitations and pre-registered experiments}
\label{sec:limitations}

\subsection{Limitations}

\begin{itemize}
    \item \textbf{The central result rests on two seeds.} In the run reported in Table~\ref{tab:curriculum}, arm~F completed on seeds 42 and 137 only, the third having diverged. The F-vs-A direction is consistent ($2/2$) but cannot be significant at $n=2$, where the exact paired permutation test floors at 0.25. A five-seed replication is in progress and is the experiment that would settle this.
    \item \textbf{F vs B is sign-split} and is not claimed (\S\ref{sec:removal}).
    \item \textbf{The hard switch is a shock that occasionally kills a run} (one divergence in six completed arm-F runs). It is stochastic rather than tied to a particular seed, and the stale-Adam-second-moment mechanism is evidenced but not yet isolated from the other two candidates.
    \item \textbf{One switch point only.} We removed windows at the 50\% mark. Whether 10\%, 25\%, or 75\% is better---or whether it matters at all---is untested, and this is what would turn the finding into an actionable recipe.
    \item \textbf{The per-layer ablation is a single seed}; the Only~$L_0,L_1$ vs Quartic margin needs replication.
    \item \textbf{The variance-reduction control (Exp 3) is a single seed}, and its batch-256 arm did not complete.
    \item \textbf{Implicit regularization is unresolved} (seed-contradictory at $n=2$).
    \item \textbf{125M evidence is a 20k-step probe}, not a converged comparison, and one of five seeds is negative. The removal experiment has not been run at 125M at all.
    \item \textbf{The endpoint measurement is noisier than the effect.} A single evaluation over 12 randomly drawn batches has a standard deviation of ${\approx}0.007$ bpb on a frozen checkpoint, against an effect of 0.0136. Paired comparison cancels most of it (residual 0.0010 within a script, 0.0023 across scripts), which is why every claim here is paired and why an unpaired test on these data would be uninformative. A fixed evaluation set would remove the term entirely; see item~(4) of the pre-registration.
    \item \textbf{Mechanism experiments were performed only at 3.4M.}
    \item \textbf{No downstream task evaluation} (CORE, MMLU, etc.); the 3.4M model uses a narrow domain (TinyStories).
    \item \textbf{The optimal exponent at depth 12 is undetermined}---only $\gamma = 4$ was run at 125M.
    \item \textbf{A second-order probe of landscape flatness} (Hessian trace and top eigenvalue) on one converged matched pair was inconclusive at $n=1$ and is not reported as a result. The attention-entropy/sharpness relationship is itself an established area; a credible test requires multiple converged pairs at scale.
\end{itemize}

\subsection{Pre-registered experiments}

We state criteria in advance so that outcomes cannot be re-framed after the fact.

\paragraph{(1) Removal at five seeds.} Run arm~F at seeds 42/137/256/789/1337 with the switch at step 10{,}000; arms A and B already exist at all five seeds. \emph{Criterion:} claim ``removal preserves the benefit'' iff all 5 paired differences (A~$-$~F) are positive, giving the exact permutation floor $p = 0.031$; report paired-$t$ and $d_z$ alongside. Claim ``removal is better than keeping'' only if at least 4 of 5 paired differences (B~$-$~F) are positive \emph{and} the paired $t$-test gives $p < 0.05$; at 2/5 or 3/5 we report no detectable difference. Diverged seeds count in the denominator.

\paragraph{(2) Shock diagnosis.} The divergence does not reproduce on demand---it occurs roughly once in six runs and is not tied to a seed---so an experiment keyed to reproducing the NaN would be badly powered. We therefore target the \emph{shock}, which is continuous, large, and reproduces on every seed. From one shared pre-switch state, run four treatments and compare peak loss, peak pre-clip gradient norm, and peak Adam update over $[\text{switch}, \text{switch}+500]$: (a) the unmodified switch, as a reference; (b) switching to an explicit all-full window list, which preserves the masked-softmax implementation and changes only the mask, separating cause~(i) from (ii) and (iii); (c) resetting the optimizer state at the switch, which removes cause~(ii) by construction; (d) a 200-step LR re-warmup and (e) a 500-step linear window ramp, as candidate mitigations. Sharing one pre-switch state makes these a controlled comparison rather than five independent runs. \emph{Criterion:} report which treatment reduces the peak Adam update toward its pre-switch value of $3.5\times10^{-3}$. Treatment~(c) is predicted to; if it does not, the stale-second-moment account in \S\ref{sec:removal} is wrong and must be withdrawn.

\paragraph{(3) Switch-point sweep.} Remove windows at 2k / 5k / 10k / 15k of a 20k-step run, at 5 seeds. The endpoints $x=0$ and $x=20\text{k}$ are arms A and B. \emph{Criterion:} claim an optimal removal point only if the best interior $x$ beats \emph{both} endpoints on at least 4 of 5 paired seeds. If the curve is flat in $x$, the finding is that the removal point does not matter over 10--75\% of training---which is a stronger recipe, since it requires no tuning. Report the divergence rate per switch point.

\paragraph{(4) Fixed evaluation set.} Every reported endpoint is a single evaluation over 12 randomly drawn batches (98{,}304 tokens, 0.51\% of the validation set), resampled on each call. Evaluating one \emph{frozen} checkpoint twelve times gives a standard deviation of 0.0083 (baseline, seed 42) and 0.0063 (quartic), with a range up to 0.030---over half the size of the 0.0136 effect. The paired design rescues this: paired arms consume the random number stream near-identically and therefore draw nearly the same evaluation batches, so the noise is common-mode. Detrended evaluation fluctuations correlate at $r=0.989$ between the two arms produced by the same script, leaving a residual difference standard deviation of 0.0010. Arm~F is produced by a different script that consumes the stream slightly differently, giving $r=0.947$ and a residual of 0.0023. Against that residual, $F$ vs $A$ is $3.9\sigma$ and $7.6\sigma$ and $B$ vs $A$ is $11$--$20\sigma$, but $F$ vs $B$ is $0.8\sigma$ and $2.1\sigma$---indistinguishable from noise, which is the quantitative reason the sign split in Table~\ref{tab:curriculum} cannot be resolved. \emph{Change:} evaluate on a deterministic set of offsets built once from a private random generator (so it draws nothing from the global stream and leaves training data order untouched), sized at ${\approx}1$M tokens. \emph{Criterion:} repeated evaluation of one frozen checkpoint must return bit-identical values, i.e.\ standard deviation exactly zero. This changes the recorded endpoint of every run, so all arms must be re-measured together rather than compared against existing numbers. \emph{Consequence for (1):} at the current residual, $n=5$ gives roughly a 22\% chance of five positive differences and an expected paired-$t$ $p\approx0.22$ for $F$ vs $B$, so experiment~(1) is not expected to resolve it; reaching $p<0.05$ at $n=5$ requires the residual below ${\approx}0.0012$.

\paragraph{(5) Converged 125M, three arms, one horizon.} Baseline / windows-throughout / windows-removed-at-25k, at 50k steps ($6.55$B tokens, ${\approx}53$ tokens per parameter, comfortably past the compute-optimal ratio of \citet{hoffmann2022chinchilla}), under a single fully-annealed cosine with no resume, at 3--5 seeds. Cost is ${\approx}5$ H100-hours per run, so ${\approx}45$ H100-hours at 3 seeds and ${\approx}75$ at 5. \emph{Criterion:} a run counts as converged only if its decline over the final 10k steps is below 0.002 bpb per 1k steps; report the measured value for every run. Note that at $n=3$ the exact permutation test floors at $0.125$ and cannot reach conventional significance, so a 3-seed result is reported descriptively with the per-seed table. \emph{We pre-register the negative outcome:} if the converged gap is smaller than the 20k gap, we will report that the 50k figures in Appendix~\ref{app:125m50k} were an undertraining artifact.

\section{Discussion}

\paragraph{The practical recipe.} If the removal result replicates, the deployment story is simple: apply a depth-wise locality constraint for roughly the first half of training, then remove it and train normally. The shipped model is a standard transformer with no architectural modification, no extra parameters, and standard inference cost. This differs from every local--global hybrid in production \citep{jiang2023mistral,gemma3} and from SWAT \citep{fu2025swat}, all of which retain windows at inference, and it means the technique composes with whatever attention implementation a deployment already uses.

\paragraph{Curriculum, not architecture.} The result that makes this a curriculum rather than an architecture is the persistence in Table~\ref{tab:entropy}: the low early-layer entropy the constraint induces survives $5\times$ longer training. The model does not revert when the mask is lifted. Combined with the gradient covariance collapse under a restricted window, the picture is that the constraint makes early-layer updates coherent during the phase when the hierarchy is being laid down, and that the resulting hierarchy is then self-sustaining. We have shown the persistence and the coherence separately; we have not shown that the second causes the first.

\paragraph{What would falsify this.} A flat or negative switch-point sweep at five seeds, or a converged 125M comparison in which the removed arm underperforms the windowed arm, would each undercut the recipe. The criteria are fixed in \S\ref{sec:limitations} before the runs.

\section{Conclusion}

Early-layer locality is a well-established prior for transformers, and we do not claim it. We claim that it is needed only during early training. Removing a quartic depth-wise window schedule at the halfway point of a 20k-step run retains the benefit over a paired full-attention baseline on both seeds that completed ($+1.05$\%, $+1.95$\%), while keeping the windows throughout gives $+1.26$\% and $+1.40$\%; whether removal is better than keeping is sign-split at $n=2$ and is not claimed. What the constraint leaves behind is measurable: early-layer attention entropy stays 52\% below baseline at $5\times$ longer training, and under a restricted window the early-layer gradient covariance collapses from effective rank 48.6 to 17.2 with 96.7\% of variance in a single component. A per-layer ablation puts the effect in layers~0--1 and shows late-layer locality hurts. At 125M the schedule gives $+2.6$\% at 20k across 5 paired seeds with one negative. The consequence, if the five-seed replication holds, is that this prior belongs in the training recipe rather than the deployed architecture: use windows as warmup, ship a standard transformer.

Code and data: \url{https://github.com/nicoveraz/neurogen}

\appendix

\section{The broader search this result came from}
\label{app:search}

The window schedule was not designed; it was the surviving candidate from an autonomous search over 200+ fixed-budget experiments following the autoresearch paradigm \citep{karpathy2026autoresearch}, in which an agent modifies training code, runs experiments, keeps improvements, and discards failures. Twenty-six approaches were organized into four categories: \textbf{(A) weight decoration}---cellular-automata initialization (grid, block-diagonal, reaction-diffusion, spectral), live CA during training, embryogenic activity-dependent CA, pre-wired circuits \citep{olsson2022incontext,anthropic2025circuits,mordvintsev2020growing,najarro2022hypernca}; \textbf{(B) architectural constraints}---the attention window schedules studied here, CA-generated attention bias; \textbf{(C) combinations}; \textbf{(D) radical modifications}---CA modulation channels, token vitality, sleep consolidation, cross-layer CA state, developmental dropout.

Table~\ref{tab:failures} records the negative results. They are reported for completeness and are not evidence for any claim in the main text.

\begin{table}[H]
\centering
\caption{Summary of unsuccessful approaches from the search.}
\label{tab:failures}
\begin{tabular}{lccl}
\toprule
Approach & Best result & Experiments & Failure mode \\
\midrule
CA initialization & $+0.6$\% & 40+ & Marginal; simpler patterns $>$ complex \\
Live CA & $-2$ to $-5$\% & 15+ & Per-step overhead kills performance \\
Embryogenic CA & $+1.3$\% & 16 & Modest; high overhead \\
Pre-wired circuits alone & $-3$ to $-5$\% & 12 & Gradients destroy pre-wired structure \\
CA modulation channels & divergence & 8 & Model collapse \\
Token vitality & divergence & 4 & Model collapse \\
Sleep consolidation & $-1$\% & 6 & Overhead $>$ benefit \\
\bottomrule
\end{tabular}
\end{table}

For reference, the validated 3.4M window variants at 20k steps across 5 matched seeds are in Table~\ref{tab:3_4m}. All five paired differences are positive for every variant, so the exact permutation test reaches its $1/32$ floor in each case; the paired $t$ and $d_z$ separate them.

\begin{table}[H]
\centering
\caption{3.4M validation results (20{,}000 steps, 5 matched seeds). $p_{\text{perm}}$ is the exact one-sided sign-flip permutation $p$ (floor $1/32$); $p_t$ is the two-sided paired $t$-test; $d_z$ is the paired effect size.}
\label{tab:3_4m}
\begin{tabular}{lccccccc}
\toprule
Config & $n$ & Mean bpb & Std & vs Baseline & $p_{\text{perm}}$ & $p_t$ & $d_z$ \\
\midrule
Baseline & 5 & 0.9002 & 0.0075 & --- & --- & --- & --- \\
Quadratic ($\gamma=2$) & 5 & 0.8911 & 0.0048 & $+1.0$\% & 0.031 & 0.012 & 1.94 \\
\textbf{Quartic ($\gamma=4$)} & 5 & \textbf{0.8866} & 0.0056 & $\mathbf{+1.5}$\% & \textbf{0.031} & \textbf{0.001} & \textbf{3.59} \\
Quad + Induction & 5 & 0.8899 & 0.0041 & $+1.1$\% & 0.031 & 0.009 & 2.09 \\
\bottomrule
\end{tabular}
\end{table}

Note that combining the schedule with pre-wired induction heads reaches only $+1.1$\%---adding the scaffold does not help beyond the constraint itself.

\section{125M extension to 50k steps (unconverged; suggestive only)}
\label{app:125m50k}

\textbf{These numbers do not support a scaling claim and are reported only for completeness.} Two of the five 125M seeds were extended to 50k steps. Neither arm satisfies the convergence criterion of \S\ref{sec:limitations}: over the final 10k steps the baselines decline 0.0098 and 0.0084 bpb per 1k steps and the quartic arms decline 0.0208 and 0.0166---all an order of magnitude above the 0.002 threshold, with the quartic arms descending faster than the baselines. A gap that widens between two curves that are both still falling is uninformative about the asymptotic gap.

\begin{table}[H]
\centering
\caption{125M gap evolution on the two seeds extended to 50k (mean of seeds 42 and 137). \textbf{Both arms unconverged at every row.}}
\label{tab:125m}
\begin{tabular}{rccc}
\toprule
Steps & Baseline bpb & Quartic bpb & Gap \\
\midrule
5k & 4.911 & 4.820 & $+1.9$\% \\
10k & 4.580 & 4.390 & $+4.2$\% \\
20k & 4.095 & 3.942 & $+3.8$\% \\
30k & 3.844 & 3.576 & $+7.0$\% \\
40k & 3.538 & 3.188 & $+9.9$\% \\
50k & 3.447 & 3.001 & $+12.9$\% \\
\bottomrule
\end{tabular}
\end{table}

Per-seed at 50k the gaps are $+17.2$\% (seed 42) and $+8.4$\% (seed 137), i.e.\ $n=2$ with a factor-of-two spread. The headline 125M number remains $+2.6$\% at 20k across 5 seeds (\S\ref{sec:scale}).

\section{3.4M extended training to 100k steps}
\label{app:100k}

To check that the advantage is not merely faster early convergence, both baseline and quartic were trained for 100k steps at seed 42 (Table~\ref{tab:100k}).

\begin{table}[H]
\centering
\caption{3.4M gap evolution over 100k steps (seed 42, $n=1$).}
\label{tab:100k}
\begin{tabular}{rcccc}
\toprule
Steps & Baseline bpb & Quartic bpb & Gap & Note \\
\midrule
1k & 1.339 & 1.326 & $+1.0$\% & \\
5k & 1.122 & 1.109 & $+1.2$\% & \\
10k & 1.051 & 1.021 & $+2.9$\% & peak gap \\
20k & 0.971 & 0.957 & $+1.4$\% & \\
50k & 0.898 & 0.893 & $+0.6$\% & \\
100k & 0.807 & 0.799 & $+1.0$\% & persists \\
\bottomrule
\end{tabular}
\end{table}

The gap peaks early ($+2.9$\% at 10k), narrows during mid-training, and does not close by 100k. The best checkpoint seen for each arm is 0.7921 (quartic, 96k) versus 0.7980 (baseline, 96k). This is a single seed, and the run-to-run noise floor on this hardware (${\approx}0.055$ bpb) exceeds the final gap, so the endpoint value should not be read as a measured effect size; the paired 20k results in Table~\ref{tab:3_4m} carry that. What this run supports is the qualitative shape---an early peak consistent with a curriculum effect, and no reversion at long horizons, which is the same conclusion the entropy measurements reach in Table~\ref{tab:entropy}.

\bibliographystyle{plainnat}

\begin{thebibliography}{99}

\bibitem[Anthropic(2025)]{anthropic2025circuits}
Anthropic.
\newblock Circuit Tracing: Revealing Computational Graphs in Language Models.
\newblock \emph{Transformer Circuits Thread}, 2025.

\bibitem[Beltagy et~al.(2020)]{beltagy2020longformer}
I.~Beltagy, M.~E. Peters, and A.~Cohan.
\newblock Longformer: The Long-Document Transformer.
\newblock \emph{arXiv:2004.05150}, 2020.

\bibitem[Cabannes et~al.(2025)]{cabannes2025short}
L.~Cabannes, M.~Beck, G.~Szilvasy, M.~Douze, M.~Lomeli, J.~Copet, P.-E. Mazar\'e, G.~Synnaeve, and H.~J\'egou.
\newblock Short Window Attention Enables Long-Term Memorization.
\newblock \emph{arXiv:2509.24552}, 2025.

\bibitem[Dao(2023)]{dao2023flashattention2}
T.~Dao.
\newblock FlashAttention-2: Faster Attention with Better Parallelism and Work Partitioning.
\newblock \emph{arXiv:2307.08691}, 2023.

\bibitem[Eldan and Li(2023)]{eldan2023tinystories}
R.~Eldan and Y.~Li.
\newblock TinyStories: How Small Can Language Models Be and Still Speak Coherent English?
\newblock \emph{arXiv:2305.07759}, 2023.

\bibitem[Fu et~al.(2025)]{fu2025swat}
Z.~Fu, W.~Song, Y.~Wang, X.~Wu, Y.~Zheng, Y.~Zhang, D.~Xu, X.~Wei, T.~Xu, and X.~Zhao.
\newblock Sliding Window Attention Training for Efficient Large Language Models.
\newblock \emph{arXiv:2502.18845}, 2025.

\bibitem[Gemma Team(2025)]{gemma3}
Gemma Team, Google DeepMind.
\newblock Gemma 3 Technical Report.
\newblock \emph{arXiv:2503.19786}, 2025.

\bibitem[Hoffmann et~al.(2022)]{hoffmann2022chinchilla}
J.~Hoffmann, S.~Borgeaud, A.~Mensch, et~al.
\newblock Training Compute-Optimal Large Language Models.
\newblock \emph{arXiv:2203.15556}, 2022.

\bibitem[Huang et~al.(2020)]{huang2020tfixup}
X.~S. Huang, F.~Perez, J.~Ba, and M.~Volkovs.
\newblock Improving Transformer Optimization Through Better Initialization.
\newblock \emph{ICML}, 2020.

\bibitem[Jiang et~al.(2023)]{jiang2023mistral}
A.~Q. Jiang, A.~Sablayrolles, A.~Mensch, C.~Bamford, D.~S. Chaplot, et~al.
\newblock Mistral 7B.
\newblock \emph{arXiv:2310.06825}, 2023.

\bibitem[Karpathy(2026)]{karpathy2026autoresearch}
A.~Karpathy.
\newblock nanochat / autoresearch.
\newblock \url{https://github.com/karpathy/autoresearch}, 2026.

\bibitem[Mordvintsev et~al.(2020)]{mordvintsev2020growing}
A.~Mordvintsev, E.~Randazzo, E.~Niklasson, and M.~Levin.
\newblock Growing Neural Cellular Automata.
\newblock \emph{Distill}, 2020.

\bibitem[Najarro and Risi(2022)]{najarro2022hypernca}
E.~Najarro and S.~Risi.
\newblock HyperNCA: Growing Neural Cellular Automata to Grow Neural Networks.
\newblock \emph{arXiv:2204.11674}, 2022.

\bibitem[Olsson et~al.(2022)]{olsson2022incontext}
C.~Olsson, N.~Elhage, N.~Nanda, et~al.
\newblock In-context Learning and Induction Heads.
\newblock \emph{Transformer Circuits Thread}, 2022.

\bibitem[Penedo et~al.(2024)]{penedo2024fineweb}
G.~Penedo, H.~Malartic, D.~Hesslow, et~al.
\newblock The FineWeb Datasets: Decanting the Web for the Finest Text Data at Scale.
\newblock \emph{arXiv:2406.17557}, 2024.

\bibitem[Press et~al.(2021)]{press2021shortformer}
O.~Press, N.~A. Smith, and M.~Lewis.
\newblock Shortformer: Better Language Modeling using Shorter Inputs.
\newblock In \emph{Proceedings of ACL}, 2021. arXiv:2012.15832.

\bibitem[Rae and Razavi(2020)]{rae2020deep}
J.~W. Rae and A.~Razavi.
\newblock Do Transformers Need Deep Long-Range Memory?
\newblock In \emph{Proceedings of ACL}, 2020. arXiv:2007.03356.

\bibitem[Sukhbaatar et~al.(2019)]{sukhbaatar2019adaptive}
S.~Sukhbaatar, E.~Grave, P.~Bojanowski, and A.~Joulin.
\newblock Adaptive Attention Span in Transformers.
\newblock In \emph{Proceedings of ACL}, 2019. arXiv:1905.07799.

\bibitem[Trockman and Kolter(2023)]{trockman2023mimetic}
A.~Trockman and J.~Z. Kolter.
\newblock Mimetic Initialization of Self-Attention Layers.
\newblock \emph{arXiv:2305.09828}, 2023.

\bibitem[Xu et~al.(2025)]{xu2025mswa}
Y.~Xu, S.~Nag, D.~Li, L.~Tian, and E.~Barsoum.
\newblock MSWA: Refining Local Attention with Multi-Scale Window Attention.
\newblock \emph{arXiv:2501.01039}, 2025.

\bibitem[Yang et~al.(2022)]{yang2022mup}
G.~Yang, E.~J. Hu, I.~Babuschkin, et~al.
\newblock Tensor Programs V: Tuning Large Neural Networks via Zero-Shot Hyperparameter Transfer.
\newblock \emph{arXiv:2203.03466}, 2022.

\end{thebibliography}

\end{document}
